\documentclass{article}
\usepackage{iclr2027_conference,times}
\renewcommand{\ttdefault}{cmtt}
%%%%% NEW MATH DEFINITIONS %%%%%

\usepackage{amsmath,amsfonts,bm}

% Mark sections of captions for referring to divisions of figures
\newcommand{\figleft}{{\em (Left)}}
\newcommand{\figcenter}{{\em (Center)}}
\newcommand{\figright}{{\em (Right)}}
\newcommand{\figtop}{{\em (Top)}}
\newcommand{\figbottom}{{\em (Bottom)}}
\newcommand{\captiona}{{\em (a)}}
\newcommand{\captionb}{{\em (b)}}
\newcommand{\captionc}{{\em (c)}}
\newcommand{\captiond}{{\em (d)}}

% Highlight a newly defined term
\newcommand{\newterm}[1]{{\bf #1}}


% Figure reference, lower-case.
\def\figref#1{figure~\ref{#1}}
% Figure reference, capital. For start of sentence
\def\Figref#1{Figure~\ref{#1}}
\def\twofigref#1#2{figures \ref{#1} and \ref{#2}}
\def\quadfigref#1#2#3#4{figures \ref{#1}, \ref{#2}, \ref{#3} and \ref{#4}}
% Section reference, lower-case.
\def\secref#1{section~\ref{#1}}
% Section reference, capital.
\def\Secref#1{Section~\ref{#1}}
% Reference to two sections.
\def\twosecrefs#1#2{sections \ref{#1} and \ref{#2}}
% Reference to three sections.
\def\secrefs#1#2#3{sections \ref{#1}, \ref{#2} and \ref{#3}}
% Reference to an equation, lower-case.
\def\eqref#1{equation~\ref{#1}}
% Reference to an equation, upper case
\def\Eqref#1{Equation~\ref{#1}}
% A raw reference to an equation---avoid using if possible
\def\plaineqref#1{\ref{#1}}
% Reference to a chapter, lower-case.
\def\chapref#1{chapter~\ref{#1}}
% Reference to an equation, upper case.
\def\Chapref#1{Chapter~\ref{#1}}
% Reference to a range of chapters
\def\rangechapref#1#2{chapters\ref{#1}--\ref{#2}}
% Reference to an algorithm, lower-case.
\def\algref#1{algorithm~\ref{#1}}
% Reference to an algorithm, upper case.
\def\Algref#1{Algorithm~\ref{#1}}
\def\twoalgref#1#2{algorithms \ref{#1} and \ref{#2}}
\def\Twoalgref#1#2{Algorithms \ref{#1} and \ref{#2}}
% Reference to a part, lower case
\def\partref#1{part~\ref{#1}}
% Reference to a part, upper case
\def\Partref#1{Part~\ref{#1}}
\def\twopartref#1#2{parts \ref{#1} and \ref{#2}}

\def\ceil#1{\lceil #1 \rceil}
\def\floor#1{\lfloor #1 \rfloor}
\def\1{\bm{1}}
\newcommand{\train}{\mathcal{D}}
\newcommand{\valid}{\mathcal{D_{\mathrm{valid}}}}
\newcommand{\test}{\mathcal{D_{\mathrm{test}}}}

\def\eps{{\epsilon}}


% Random variables
\def\reta{{\textnormal{$\eta$}}}
\def\ra{{\textnormal{a}}}
\def\rb{{\textnormal{b}}}
\def\rc{{\textnormal{c}}}
\def\rd{{\textnormal{d}}}
\def\re{{\textnormal{e}}}
\def\rf{{\textnormal{f}}}
\def\rg{{\textnormal{g}}}
\def\rh{{\textnormal{h}}}
\def\ri{{\textnormal{i}}}
\def\rj{{\textnormal{j}}}
\def\rk{{\textnormal{k}}}
\def\rl{{\textnormal{l}}}
% rm is already a command, just don't name any random variables m
\def\rn{{\textnormal{n}}}
\def\ro{{\textnormal{o}}}
\def\rp{{\textnormal{p}}}
\def\rq{{\textnormal{q}}}
\def\rr{{\textnormal{r}}}
\def\rs{{\textnormal{s}}}
\def\rt{{\textnormal{t}}}
\def\ru{{\textnormal{u}}}
\def\rv{{\textnormal{v}}}
\def\rw{{\textnormal{w}}}
\def\rx{{\textnormal{x}}}
\def\ry{{\textnormal{y}}}
\def\rz{{\textnormal{z}}}

% Random vectors
\def\rvepsilon{{\mathbf{\epsilon}}}
\def\rvtheta{{\mathbf{\theta}}}
\def\rva{{\mathbf{a}}}
\def\rvb{{\mathbf{b}}}
\def\rvc{{\mathbf{c}}}
\def\rvd{{\mathbf{d}}}
\def\rve{{\mathbf{e}}}
\def\rvf{{\mathbf{f}}}
\def\rvg{{\mathbf{g}}}
\def\rvh{{\mathbf{h}}}
\def\rvu{{\mathbf{i}}}
\def\rvj{{\mathbf{j}}}
\def\rvk{{\mathbf{k}}}
\def\rvl{{\mathbf{l}}}
\def\rvm{{\mathbf{m}}}
\def\rvn{{\mathbf{n}}}
\def\rvo{{\mathbf{o}}}
\def\rvp{{\mathbf{p}}}
\def\rvq{{\mathbf{q}}}
\def\rvr{{\mathbf{r}}}
\def\rvs{{\mathbf{s}}}
\def\rvt{{\mathbf{t}}}
\def\rvu{{\mathbf{u}}}
\def\rvv{{\mathbf{v}}}
\def\rvw{{\mathbf{w}}}
\def\rvx{{\mathbf{x}}}
\def\rvy{{\mathbf{y}}}
\def\rvz{{\mathbf{z}}}

% Elements of random vectors
\def\erva{{\textnormal{a}}}
\def\ervb{{\textnormal{b}}}
\def\ervc{{\textnormal{c}}}
\def\ervd{{\textnormal{d}}}
\def\erve{{\textnormal{e}}}
\def\ervf{{\textnormal{f}}}
\def\ervg{{\textnormal{g}}}
\def\ervh{{\textnormal{h}}}
\def\ervi{{\textnormal{i}}}
\def\ervj{{\textnormal{j}}}
\def\ervk{{\textnormal{k}}}
\def\ervl{{\textnormal{l}}}
\def\ervm{{\textnormal{m}}}
\def\ervn{{\textnormal{n}}}
\def\ervo{{\textnormal{o}}}
\def\ervp{{\textnormal{p}}}
\def\ervq{{\textnormal{q}}}
\def\ervr{{\textnormal{r}}}
\def\ervs{{\textnormal{s}}}
\def\ervt{{\textnormal{t}}}
\def\ervu{{\textnormal{u}}}
\def\ervv{{\textnormal{v}}}
\def\ervw{{\textnormal{w}}}
\def\ervx{{\textnormal{x}}}
\def\ervy{{\textnormal{y}}}
\def\ervz{{\textnormal{z}}}

% Random matrices
\def\rmA{{\mathbf{A}}}
\def\rmB{{\mathbf{B}}}
\def\rmC{{\mathbf{C}}}
\def\rmD{{\mathbf{D}}}
\def\rmE{{\mathbf{E}}}
\def\rmF{{\mathbf{F}}}
\def\rmG{{\mathbf{G}}}
\def\rmH{{\mathbf{H}}}
\def\rmI{{\mathbf{I}}}
\def\rmJ{{\mathbf{J}}}
\def\rmK{{\mathbf{K}}}
\def\rmL{{\mathbf{L}}}
\def\rmM{{\mathbf{M}}}
\def\rmN{{\mathbf{N}}}
\def\rmO{{\mathbf{O}}}
\def\rmP{{\mathbf{P}}}
\def\rmQ{{\mathbf{Q}}}
\def\rmR{{\mathbf{R}}}
\def\rmS{{\mathbf{S}}}
\def\rmT{{\mathbf{T}}}
\def\rmU{{\mathbf{U}}}
\def\rmV{{\mathbf{V}}}
\def\rmW{{\mathbf{W}}}
\def\rmX{{\mathbf{X}}}
\def\rmY{{\mathbf{Y}}}
\def\rmZ{{\mathbf{Z}}}

% Elements of random matrices
\def\ermA{{\textnormal{A}}}
\def\ermB{{\textnormal{B}}}
\def\ermC{{\textnormal{C}}}
\def\ermD{{\textnormal{D}}}
\def\ermE{{\textnormal{E}}}
\def\ermF{{\textnormal{F}}}
\def\ermG{{\textnormal{G}}}
\def\ermH{{\textnormal{H}}}
\def\ermI{{\textnormal{I}}}
\def\ermJ{{\textnormal{J}}}
\def\ermK{{\textnormal{K}}}
\def\ermL{{\textnormal{L}}}
\def\ermM{{\textnormal{M}}}
\def\ermN{{\textnormal{N}}}
\def\ermO{{\textnormal{O}}}
\def\ermP{{\textnormal{P}}}
\def\ermQ{{\textnormal{Q}}}
\def\ermR{{\textnormal{R}}}
\def\ermS{{\textnormal{S}}}
\def\ermT{{\textnormal{T}}}
\def\ermU{{\textnormal{U}}}
\def\ermV{{\textnormal{V}}}
\def\ermW{{\textnormal{W}}}
\def\ermX{{\textnormal{X}}}
\def\ermY{{\textnormal{Y}}}
\def\ermZ{{\textnormal{Z}}}

% Vectors
\def\vzero{{\bm{0}}}
\def\vone{{\bm{1}}}
\def\vmu{{\bm{\mu}}}
\def\vtheta{{\bm{\theta}}}
\def\va{{\bm{a}}}
\def\vb{{\bm{b}}}
\def\vc{{\bm{c}}}
\def\vd{{\bm{d}}}
\def\ve{{\bm{e}}}
\def\vf{{\bm{f}}}
\def\vg{{\bm{g}}}
\def\vh{{\bm{h}}}
\def\vi{{\bm{i}}}
\def\vj{{\bm{j}}}
\def\vk{{\bm{k}}}
\def\vl{{\bm{l}}}
\def\vm{{\bm{m}}}
\def\vn{{\bm{n}}}
\def\vo{{\bm{o}}}
\def\vp{{\bm{p}}}
\def\vq{{\bm{q}}}
\def\vr{{\bm{r}}}
\def\vs{{\bm{s}}}
\def\vt{{\bm{t}}}
\def\vu{{\bm{u}}}
\def\vv{{\bm{v}}}
\def\vw{{\bm{w}}}
\def\vx{{\bm{x}}}
\def\vy{{\bm{y}}}
\def\vz{{\bm{z}}}

% Elements of vectors
\def\evalpha{{\alpha}}
\def\evbeta{{\beta}}
\def\evepsilon{{\epsilon}}
\def\evlambda{{\lambda}}
\def\evomega{{\omega}}
\def\evmu{{\mu}}
\def\evpsi{{\psi}}
\def\evsigma{{\sigma}}
\def\evtheta{{\theta}}
\def\eva{{a}}
\def\evb{{b}}
\def\evc{{c}}
\def\evd{{d}}
\def\eve{{e}}
\def\evf{{f}}
\def\evg{{g}}
\def\evh{{h}}
\def\evi{{i}}
\def\evj{{j}}
\def\evk{{k}}
\def\evl{{l}}
\def\evm{{m}}
\def\evn{{n}}
\def\evo{{o}}
\def\evp{{p}}
\def\evq{{q}}
\def\evr{{r}}
\def\evs{{s}}
\def\evt{{t}}
\def\evu{{u}}
\def\evv{{v}}
\def\evw{{w}}
\def\evx{{x}}
\def\evy{{y}}
\def\evz{{z}}

% Matrix
\def\mA{{\bm{A}}}
\def\mB{{\bm{B}}}
\def\mC{{\bm{C}}}
\def\mD{{\bm{D}}}
\def\mE{{\bm{E}}}
\def\mF{{\bm{F}}}
\def\mG{{\bm{G}}}
\def\mH{{\bm{H}}}
\def\mI{{\bm{I}}}
\def\mJ{{\bm{J}}}
\def\mK{{\bm{K}}}
\def\mL{{\bm{L}}}
\def\mM{{\bm{M}}}
\def\mN{{\bm{N}}}
\def\mO{{\bm{O}}}
\def\mP{{\bm{P}}}
\def\mQ{{\bm{Q}}}
\def\mR{{\bm{R}}}
\def\mS{{\bm{S}}}
\def\mT{{\bm{T}}}
\def\mU{{\bm{U}}}
\def\mV{{\bm{V}}}
\def\mW{{\bm{W}}}
\def\mX{{\bm{X}}}
\def\mY{{\bm{Y}}}
\def\mZ{{\bm{Z}}}
\def\mBeta{{\bm{\beta}}}
\def\mPhi{{\bm{\Phi}}}
\def\mLambda{{\bm{\Lambda}}}
\def\mSigma{{\bm{\Sigma}}}

% Tensor
\DeclareMathAlphabet{\mathsfit}{\encodingdefault}{\sfdefault}{m}{sl}
\SetMathAlphabet{\mathsfit}{bold}{\encodingdefault}{\sfdefault}{bx}{n}
\newcommand{\tens}[1]{\bm{\mathsfit{#1}}}
\def\tA{{\tens{A}}}
\def\tB{{\tens{B}}}
\def\tC{{\tens{C}}}
\def\tD{{\tens{D}}}
\def\tE{{\tens{E}}}
\def\tF{{\tens{F}}}
\def\tG{{\tens{G}}}
\def\tH{{\tens{H}}}
\def\tI{{\tens{I}}}
\def\tJ{{\tens{J}}}
\def\tK{{\tens{K}}}
\def\tL{{\tens{L}}}
\def\tM{{\tens{M}}}
\def\tN{{\tens{N}}}
\def\tO{{\tens{O}}}
\def\tP{{\tens{P}}}
\def\tQ{{\tens{Q}}}
\def\tR{{\tens{R}}}
\def\tS{{\tens{S}}}
\def\tT{{\tens{T}}}
\def\tU{{\tens{U}}}
\def\tV{{\tens{V}}}
\def\tW{{\tens{W}}}
\def\tX{{\tens{X}}}
\def\tY{{\tens{Y}}}
\def\tZ{{\tens{Z}}}


% Graph
\def\gA{{\mathcal{A}}}
\def\gB{{\mathcal{B}}}
\def\gC{{\mathcal{C}}}
\def\gD{{\mathcal{D}}}
\def\gE{{\mathcal{E}}}
\def\gF{{\mathcal{F}}}
\def\gG{{\mathcal{G}}}
\def\gH{{\mathcal{H}}}
\def\gI{{\mathcal{I}}}
\def\gJ{{\mathcal{J}}}
\def\gK{{\mathcal{K}}}
\def\gL{{\mathcal{L}}}
\def\gM{{\mathcal{M}}}
\def\gN{{\mathcal{N}}}
\def\gO{{\mathcal{O}}}
\def\gP{{\mathcal{P}}}
\def\gQ{{\mathcal{Q}}}
\def\gR{{\mathcal{R}}}
\def\gS{{\mathcal{S}}}
\def\gT{{\mathcal{T}}}
\def\gU{{\mathcal{U}}}
\def\gV{{\mathcal{V}}}
\def\gW{{\mathcal{W}}}
\def\gX{{\mathcal{X}}}
\def\gY{{\mathcal{Y}}}
\def\gZ{{\mathcal{Z}}}

% Sets
\def\sA{{\mathbb{A}}}
\def\sB{{\mathbb{B}}}
\def\sC{{\mathbb{C}}}
\def\sD{{\mathbb{D}}}
% Don't use a set called E, because this would be the same as our symbol
% for expectation.
\def\sF{{\mathbb{F}}}
\def\sG{{\mathbb{G}}}
\def\sH{{\mathbb{H}}}
\def\sI{{\mathbb{I}}}
\def\sJ{{\mathbb{J}}}
\def\sK{{\mathbb{K}}}
\def\sL{{\mathbb{L}}}
\def\sM{{\mathbb{M}}}
\def\sN{{\mathbb{N}}}
\def\sO{{\mathbb{O}}}
\def\sP{{\mathbb{P}}}
\def\sQ{{\mathbb{Q}}}
\def\sR{{\mathbb{R}}}
\def\sS{{\mathbb{S}}}
\def\sT{{\mathbb{T}}}
\def\sU{{\mathbb{U}}}
\def\sV{{\mathbb{V}}}
\def\sW{{\mathbb{W}}}
\def\sX{{\mathbb{X}}}
\def\sY{{\mathbb{Y}}}
\def\sZ{{\mathbb{Z}}}

% Entries of a matrix
\def\emLambda{{\Lambda}}
\def\emA{{A}}
\def\emB{{B}}
\def\emC{{C}}
\def\emD{{D}}
\def\emE{{E}}
\def\emF{{F}}
\def\emG{{G}}
\def\emH{{H}}
\def\emI{{I}}
\def\emJ{{J}}
\def\emK{{K}}
\def\emL{{L}}
\def\emM{{M}}
\def\emN{{N}}
\def\emO{{O}}
\def\emP{{P}}
\def\emQ{{Q}}
\def\emR{{R}}
\def\emS{{S}}
\def\emT{{T}}
\def\emU{{U}}
\def\emV{{V}}
\def\emW{{W}}
\def\emX{{X}}
\def\emY{{Y}}
\def\emZ{{Z}}
\def\emSigma{{\Sigma}}

% entries of a tensor
% Same font as tensor, without \bm wrapper
\newcommand{\etens}[1]{\mathsfit{#1}}
\def\etLambda{{\etens{\Lambda}}}
\def\etA{{\etens{A}}}
\def\etB{{\etens{B}}}
\def\etC{{\etens{C}}}
\def\etD{{\etens{D}}}
\def\etE{{\etens{E}}}
\def\etF{{\etens{F}}}
\def\etG{{\etens{G}}}
\def\etH{{\etens{H}}}
\def\etI{{\etens{I}}}
\def\etJ{{\etens{J}}}
\def\etK{{\etens{K}}}
\def\etL{{\etens{L}}}
\def\etM{{\etens{M}}}
\def\etN{{\etens{N}}}
\def\etO{{\etens{O}}}
\def\etP{{\etens{P}}}
\def\etQ{{\etens{Q}}}
\def\etR{{\etens{R}}}
\def\etS{{\etens{S}}}
\def\etT{{\etens{T}}}
\def\etU{{\etens{U}}}
\def\etV{{\etens{V}}}
\def\etW{{\etens{W}}}
\def\etX{{\etens{X}}}
\def\etY{{\etens{Y}}}
\def\etZ{{\etens{Z}}}

% The true underlying data generating distribution
\newcommand{\pdata}{p_{\rm{data}}}
% The empirical distribution defined by the training set
\newcommand{\ptrain}{\hat{p}_{\rm{data}}}
\newcommand{\Ptrain}{\hat{P}_{\rm{data}}}
% The model distribution
\newcommand{\pmodel}{p_{\rm{model}}}
\newcommand{\Pmodel}{P_{\rm{model}}}
\newcommand{\ptildemodel}{\tilde{p}_{\rm{model}}}
% Stochastic autoencoder distributions
\newcommand{\pencode}{p_{\rm{encoder}}}
\newcommand{\pdecode}{p_{\rm{decoder}}}
\newcommand{\precons}{p_{\rm{reconstruct}}}

\newcommand{\laplace}{\mathrm{Laplace}} % Laplace distribution

\newcommand{\E}{\mathbb{E}}
\newcommand{\Ls}{\mathcal{L}}
\newcommand{\R}{\mathbb{R}}
\newcommand{\emp}{\tilde{p}}
\newcommand{\lr}{\alpha}
\newcommand{\reg}{\lambda}
\newcommand{\rect}{\mathrm{rectifier}}
\newcommand{\softmax}{\mathrm{softmax}}
\newcommand{\sigmoid}{\sigma}
\newcommand{\softplus}{\zeta}
\newcommand{\KL}{D_{\mathrm{KL}}}
\newcommand{\Var}{\mathrm{Var}}
\newcommand{\standarderror}{\mathrm{SE}}
\newcommand{\Cov}{\mathrm{Cov}}
% Wolfram Mathworld says $L^2$ is for function spaces and $\ell^2$ is for vectors
% But then they seem to use $L^2$ for vectors throughout the site, and so does
% wikipedia.
\newcommand{\normlzero}{L^0}
\newcommand{\normlone}{L^1}
\newcommand{\normltwo}{L^2}
\newcommand{\normlp}{L^p}
\newcommand{\normmax}{L^\infty}

\newcommand{\parents}{Pa} % See usage in notation.tex. Chosen to match Daphne's book.

\DeclareMathOperator*{\argmax}{arg\,max}
\DeclareMathOperator*{\argmin}{arg\,min}

\DeclareMathOperator{\sign}{sign}
\DeclareMathOperator{\Tr}{Tr}
\let\ab\allowbreak


\usepackage{hyperref}
\usepackage{url}
\usepackage{booktabs}
\usepackage{tabularx}
\usepackage{array}
\usepackage{graphicx}
\usepackage{float}
\usepackage{subcaption}
\usepackage{amsmath,amssymb,amsfonts,amsthm}
\usepackage{bm}
\usepackage{xcolor}
\usepackage{listings}

\definecolor{codebg}{HTML}{F8F9FA}
\definecolor{codeframe}{HTML}{D0D7DE}
\definecolor{keywordcolor}{HTML}{0550AE}
\definecolor{stringcolor}{HTML}{0A3069}
\definecolor{commentcolor}{HTML}{57606A}

\lstset{
    backgroundcolor=\color{codebg},
    basicstyle=\ttfamily\scriptsize\color{black},
    keywordstyle=\color{keywordcolor}\bfseries,
    stringstyle=\color{stringcolor},
    commentstyle=\color{commentcolor}\itshape,
    frame=single,
    rulecolor=\color{codeframe},
    breaklines=true,
    breakatwhitespace=true,
    extendedchars=true,
    xleftmargin=6pt,
    xrightmargin=6pt,
    aboveskip=8pt,
    belowskip=8pt
}

\hypersetup{
    colorlinks=true,
    linkcolor=blue!70!black,
    citecolor=blue!70!black,
    urlcolor=blue!70!black
}

% Metadata for salt
\title{\textbf{SALT: Singular Value Adaptation with Low-Rank Tuning for Parameter-Efficient Attention}}

\author{
  \textbf{Taha Majs}\thanks{Corresponding author: tahamajs@uni.edu} \quad \textbf{Research Team} \\
  Department of Computer Science \\
  Advanced Machine Learning \& Efficiency Laboratory
}

\date{\today}


\newcommand{\fix}{\marginpar{FIX}}
\newcommand{\new}{\marginpar{NEW}}

\begin{document}

\maketitle

\begin{abstract}
Low-rank adaptation methods such as LoRA freeze pretrained weights and add a 
trainable rank-$r$ residual. While LoRA achieves large reductions in trainable 
parameters, it does not reduce \textbf{total} parameters -- the base weight matrix 
remains at full size. We introduce \textbf{SALT} (\textbf{S}ingular Value 
\textbf{A}daptation with \textbf{L}ow-Rank \textbf{T}uning), which replaces each weight matrix 
$W \in \mathbb{R}^{d_\text{out} \times d_\text{in}}$ with a three-component 
representation: (i) a truncated SVD $(U_{:,:r}, \Sigma_{:r}, V_{:r,:})$ stored 
as \textbf{non-trainable buffers}; (ii) a trainable per-singular-value scale 
$\gamma \in \mathbb{R}^r$ that modulates the retained singular values; and (iii) 
a trainable low-rank residual $AB$ initialized to zero.

For $d_\text{out} = d_\text{in} = 768$ and $r = 16$, this reduces stored 
parameters from 589,824 to \textbf{49,936} -- a \textbf{11.8$$\times$ reduction} -- while 
maintaining the full expressiveness of a rank-$2r$ effective weight matrix. 
We evaluate on DistilBERT, RoBERTa, and GPT-2 on SST-2. Accuracy is preserved 
within \textbf{-10.43\%} (DistilBERT), \textbf{-27.64\%} (RoBERTa), and 
\textbf{-28.44\%} (GPT-2) after 50 recovery steps.

We honestly report a key limitation: because our forward pass reconstructs 
the full weight matrix $W = U \text{diag}(\Sigma \odot \gamma) V^\top + AB$ 
before calling \texttt{F.linear}, \textbf{inference memory and latency are not reduced}. 
Storage footprint is reduced by \textbf{32$$\times$}, but runtime 
efficiency is unchanged (or slightly worse due to reconstruction overhead). 
We identify the necessary change -- performing the forward as three sequential 
matmuls -- to realize the inference-time benefits.

\textbf{Keywords:} singular value decomposition, low-rank adaptation, parameter 
efficiency, fine-tuning, transformer compression



\vspace{0.5em}
\noindent\textbf{Keywords:} singular value decomposition, low-rank adaptation, parameter
\end{abstract}


\section{Introduction}

\subsection{Motivation}

Parameter-efficient fine-tuning (PEFT) has become the default approach for 
adapting large language models. Among PEFT methods, LoRA (Hu et al., 2021) 
is the most widely used: it freezes the base weight $W$ and adds a trainable 
rank-$r$ update:

\[
W_\text{LoRA} = W_\text{frozen} + \frac{\alpha}{r} B A
\tag{1.1}
\]

\textbf{Limitation.} LoRA reduces \textbf{trainable} parameters, but not \textbf{total} 
parameters. The base weight $W_\text{frozen}$ remains at full size. For a 
7B-parameter model, LoRA reduces the trainable count from 7B to ~10M but 
the checkpoint still weighs 7B parameters.

\textbf{Alternative: low-rank weights.} If $W$ itself is low-rank, we can store 
it as $W = U V$ with $r \ll \min(d_\text{in}, d_\text{out})$. This reduces 
total parameters. But it also throws away information: the base model no 
longer matches its pretrained state, and downstream accuracy may collapse.

0.85


\subsection{Problem Statement}

We seek a parameterization $\mathcal{P}(W)$ such that:

\begin{enumerate}
\item \textbf{Total parameter reduction:} $|\mathcal{P}(W)| \ll |W|$
\item \textbf{Expressiveness:} $\mathcal{P}$ can represent a rich family of matrices,
  not just rank-$r$ ones
\item \textbf{Pretrained initialization:} $\mathcal{P}$ can be initialized to closely
  match a given $W_\text{pretrained}$
\item \textbf{Trainability:} a small subset of $\mathcal{P}$ is trainable

\end{enumerate}
\subsection{Our Approach}

SALT parameterizes $W$ as:

\[
W = U_{:,:r} \cdot \text{diag}(\Sigma_{:r} \odot \gamma) \cdot V_{:r,:}^\top + A B
\tag{1.2}
\]

where:

\begin{itemize}
\item $U_{:,:r} \in \mathbb{R}^{d_\text{out} \times r}$,
  $\Sigma_{:r} \in \mathbb{R}^r$,
  $V_{:r,:} \in \mathbb{R}^{r \times d_\text{in}}$ -- \textbf{buffers} from the SVD of $W_\text{pretrained}$
\item $\gamma \in \mathbb{R}^r$ -- trainable per-singular-value scale (init: 1)
\item $A \in \mathbb{R}^{d_\text{out} \times r}$,
  $B \in \mathbb{R}^{r \times d_\text{in}}$ -- trainable low-rank residual (init: A small random, B = 0)

\end{itemize}
\textbf{At initialization:}
\begin{itemize}
\item $\gamma = \mathbf{1} \Rightarrow U \text{diag}(\Sigma) V^\top$ = truncated SVD of $W_\text{pretrained}$
\item $A B = 0 \Rightarrow$ no residual contribution

\end{itemize}
So SALT starts as a rank-$r$ approximation of the pretrained weight.

\subsection{Contributions}

\begin{itemize}
\item \textbf{A new parameterization} that combines truncated SVD with trainable
  singular-value modulation and a low-rank residual.
\item \textbf{A 12$\times$ parameter reduction} for 768$\times$768 matrices at rank 16.
\item \textbf{A complete implementation} supporting four transformer families with
  multi-architecture weight extraction.
\item \textbf{An honest evaluation} showing that SALT preserves accuracy within
  \textbf{0.85\%} but that our naive forward implementation does not
  realize inference-time benefits.

\end{itemize}
\subsection{Paper Organization}

Section 2 reviews low-rank adaptation and factorization methods. Section 3 
provides background on SVD, LoRA, and rank-$r$ approximation. Section 4 
presents SALT. Section 5 describes integration. Section 6 details the setup. 
Section 7 reports results. Section 8 analyzes the trade-off between storage 
and compute. Section 9 discusses guidance. Section 10 lists limitations. 
Section 11 concludes.



\subsection{Notation}

\begin{table}[H]
\centering
\small
\begin{tabular}{lll}
\toprule
\textbf{Symbol} & \textbf{Meaning} & \textbf{First appears} \\
\midrule
$W$ \& Weight matrix & 1.1 \\
$d_\text{in}, d_\text{out}$ \& Input/output dims & 1.1 \\
$r$ \& SALT rank & 1.1 \\
$U, \Sigma, V$ \& SVD factors & 3.1 \\
$\gamma$ \& Singular value scale & 1.3 \\
$A, B$ \& Low-rank residual & 1.3 \\
$\alpha$ \& LoRA scaling & 3.4 \\
\bottomrule
\end{tabular}
\end{table}
</parameter>
\begin{lstlisting}

---
\end{lstlisting}

\section{Related Work}

\subsection{Low-Rank Adaptation}

\subsubsection{LoRA}

LoRA (Hu et al., 2021) parameterizes updates as $\Delta W = \frac{\alpha}{r} BA$. 
Base weights are frozen.

\textbf{Limitation:} Total parameters unchanged.

\subsubsection{AdaLoRA}

AdaLoRA (Zhang et al., 2023) adaptively allocates rank across layers.

\textbf{Limitation:} Still freezes full base weights.

\subsubsection{DoRA}

DoRA (Liu et al., 2024) decomposes weights into magnitude and direction.

\textbf{Limitation:} Magnitude is scalar per output; still stores full W.

\subsection{Low-Rank Weights}

\subsubsection{SVD Compression}

Classical SVD compression (Denton et al., 2014) replaces W with $U_r \Sigma_r V_r^\top$.

\textbf{Limitation:} No adaptation after compression; accuracy may drop.

\subsubsection{Tensor Decomposition}

Tensor train (Novikov et al., 2015) and Tucker decomposition.

\textbf{Limitation:} Complex to train; not pretrained-friendly.

\subsubsection{ALBERT Factorized Embeddings}

ALBERT (Lan et al., 2020) factorizes the embedding matrix.

\textbf{Limitation:} Only embeddings, not attention weights.

\subsection{Hybrid Methods}

\subsubsection{PiSSA}

PiSSA (Meng et al., 2024) initializes LoRA factors from the SVD of W.

\textbf{Difference from SALT:} PiSSA still stores the full W; SALT stores only 
$U_r, \Sigma_r, V_r$.

\subsubsection{OLoRA}

OLoRA (B y kaky z, 2024) initializes LoRA with orthogonal matrices.

\textbf{Difference:} OLoRA does not reduce total parameters.

\subsubsection{GaLore}

GaLore (Zhao et al., 2024) projects gradients into a low-rank subspace.

\textbf{Difference:} GaLore reduces gradient memory, not weight storage.

\subsection{Position of Our Work}

\begin{table}[H]
\centering
\small
\begin{tabular}{lllll}
\toprule
\textbf{Axis} & \textbf{LoRA} & \textbf{PiSSA} & \textbf{SVD Comp} & \textbf{SALT} \\
\midrule
Total params reduced? \& No \& No \& Yes & \textbf{Yes} \\
Trainable adaptation? \& Yes \& Yes \& No & \textbf{Yes} \\
Pretrained init? \& Base only \& SVD init \& Base & \textbf{SVD + $\gamma$ + AB} \\
Inference speedup? \& No \& No \& Yes (if fused) \& No (naive) \\
\bottomrule
\end{tabular}
\end{table}

\textbf{Key distinction.} SALT reduces \textbf{total} parameters AND supports 
adaptation AND starts from the pretrained state. Prior work achieves at most 
two of these three.



\subsection{Comparison Table}

\begin{lstlisting}

---
\end{lstlisting}

\section{Background}

\subsection{Singular Value Decomposition}

For $W \in \mathbb{R}^{m \times n}$:

\[
W = U \Sigma V^\top
\tag{3.1}
\]

where $U \in \mathbb{R}^{m \times m}$, $V \in \mathbb{R}^{n \times n}$ are 
orthogonal, and $\Sigma$ is diagonal with $\sigma_1 \geq \sigma_2 \geq \dots \geq 0$.

\textbf{Truncated SVD.}

\[
W_r = U_{:, :r} \Sigma_{:r} V_{:r, :}^\top
\tag{3.2}
\]

This is the best rank-$r$ approximation in Frobenius norm (Eckart-Young).

\subsection{Low-Rank Adaptation (LoRA)}

LoRA:

\[
W_\text{eff} = W + \frac{\alpha}{r} BA
\tag{3.3}
\]

with $A \in \mathbb{R}^{r \times d_\text{in}}$, $B \in \mathbb{R}^{d_\text{out} \times r}$.

\subsection{Singular Value Modulation}

Given fixed $U, \Sigma, V$, one can modulate the singular values via a 
diagonal scale:

\[
W(\gamma) = U \text{diag}(\Sigma \odot \gamma) V^\top
\tag{3.4}
\]

with $\gamma \in \mathbb{R}^r$. Setting $\gamma = \mathbf{1}$ recovers the 
original.

\textbf{Effect.} Increasing $\gamma_i$ amplifies the $i$-th singular direction. 
Decreasing suppresses it.

\subsection{Parameter Counting}

For a weight $W \in \mathbb{R}^{m \times n}$ at rank $r$:

\begin{table}[H]
\centering
\small
\begin{tabular}{ll}
\toprule
\textbf{Representation} & \textbf{Parameters} \\
\midrule
Full W & $mn$ \\
Truncated SVD ($U_r, \Sigma_r, V_r$) & $r(m + n + 1)$ \\
LoRA ($A, B$) & $r(m + n)$ \\
SALT ($U_r, \Sigma_r, V_r, \gamma, A, B$) & $2r(m + n) + 2r$ \\
\bottomrule
\end{tabular}
\end{table}

\textbf{Example (m = n = 768, r = 16):}

\begin{itemize}
\item Full: 589,824
\item Truncated SVD: 24,592
\item LoRA: 24,576
\item SALT: 49,184 $\to$ \textbf{~12$$\times$ reduction vs full}

\end{itemize}
\subsection{Forward Pass Efficiency}

At inference, computing $y = Wx$ requires:

\begin{itemize}
\item \textbf{Full W:} one matmul $Wx$ -- cost $\mathcal{O}(mn)$
\item \textbf{Low-rank W = UV:} two matmuls $Vx$ then $U(\cdot)$ -- cost $\mathcal{O}(r(m+n))$
\item \textbf{SALT (naive):} reconstruct W then one matmul -- cost $\mathcal{O}(mn + r(m+n))$, \textbf{slower}
\item \textbf{SALT (fused):} three matmuls -- cost $\mathcal{O}(r(m+n))$

\end{itemize}
\textbf{Observation.} Realizing the inference benefit requires \textbf{fused kernels} 
or a forward pass that avoids materializing W.

\subsection{Notation Summary}

\begin{table}[H]
\centering
\small
\begin{tabular}{ll}
\toprule
\textbf{Symbol} & \textbf{Meaning} \\
\midrule
$W$ \& Weight matrix \\
$U, \Sigma, V$ \& SVD factors \\
$r$ \& Rank \\
$\gamma$ \& Singular value scale \\
$A, B$ \& Low-rank residual \\
$d_\text{in}, d_\text{out}$ \& Input/output dims \\
\bottomrule
\end{tabular}
\end{table}
</parameter>
\begin{lstlisting}

---
\end{lstlisting}

\section{Method}

\subsection{Overview}

SALT replaces each linear projection $W$ with a three-component parameterization:

\begin{lstlisting}
    Weight W $\in$ R^{d_out $$\times$ d_in}
            |
             
    +-------------------------+
    |  Truncated SVD          |
    |  W $\approx$ U_r $\Sigma$_r V_r^T      |
    +-----------+-------------+
                |
    +-----------+------------+
                              
    +--------------+   +-----------------+
    |  Buffers     |   |  Trainable      |
    |  U, S, Vt    |   |  $\gamma$, A, B        |
    |  (frozen)    |   |                 |
    +--------------+   +-----------------+
                |                 |
                +--------+--------+
                          
    +--------------------------------------+
    |  Forward:                             |
    |  W_eff = U diag(S   $\gamma$) V^T + A B     |
    |  y = W_eff x + b                      |
    +--------------------------------------+
\end{lstlisting}

\subsection{The SVD Skeleton}

Given a pretrained weight $W \in \mathbb{R}^{d_\text{out} \times d_\text{in}}$, 
compute its SVD:

\[
W = U \Sigma V^\top, \quad U \in \mathbb{R}^{d_\text{out} \times d_\text{out}}, 
\Sigma \in \mathbb{R}^{d_\text{out}}, V \in \mathbb{R}^{d_\text{in} \times d_\text{in}}
\tag{4.1}
\]

Keep the top-$r$:

\[
U_r = U_{:, :r}, \quad S_r = \Sigma_{:r}, \quad V_r = V_{:r, :}^\top
\tag{4.2}
\]

We store $U_r, S_r, V_r$ as \textbf{buffers} (non-trainable), typically in the 
model's dtype (FP16 or FP32).

\subsection{Singular Value Modulation}

We add a trainable scale $\gamma \in \mathbb{R}^r$, initialized to 1:

\[
\Sigma_\gamma = S_r \odot \gamma
\tag{4.3}
\]

This allows the model to amplify or suppress retained singular directions.

\textbf{Why is this useful?}
\begin{itemize}
\item The pretrained SVD is optimal in Frobenius norm but not for the downstream
  task.
\item $\gamma$ provides a small number of trainable parameters ($r$) that adjust
  the spectral profile.

\end{itemize}
\subsection{Low-Rank Residual}

We add a trainable low-rank residual $AB$:

\[
A \in \mathbb{R}^{d_\text{out} \times r}, \quad 
B \in \mathbb{R}^{r \times d_\text{in}}
\tag{4.4}
\]

\textbf{Initialization:}
\begin{itemize}
\item $A \sim \mathcal{N}(0, 0.01^2)$ (small random)
\item $B = 0$

\end{itemize}
So $AB = 0$ at initialization.

\textbf{Why a residual?} The SVD skeleton captures the top-$r$ directions. If 
these are insufficient, the residual can add capacity -- but starting from 
zero avoids disrupting the pretrained state.

\subsection{Combined Weight}

The effective weight is:

\[
W_\text{SALT} = U_r \cdot \text{diag}(S_r \odot \gamma) \cdot V_r + A B
\tag{4.5}
\]

\textbf{At initialization:} $W_\text{SALT} = U_r \text{diag}(S_r) V_r$ -- the 
rank-$r$ truncated SVD of $W_\text{pretrained}$.

\subsection{Forward Pass}

\subsubsection{Naive (Our Implementation)}

\begin{lstlisting}
def forward(self, x):
    adapted_S = self.S * self.gamma
    W_svd = (self.U * adapted_S.unsqueeze(0)) @ self.Vt  # reconstruct
    W_lora = self.A @ self.B                             # reconstruct
    W = (W_svd + W_lora).to(dtype=x.dtype)
    return F.linear(x, W, self.bias)                     # full matmul
\end{lstlisting}

\textbf{Cost.} $\mathcal{O}(d_\text{out} \cdot r \cdot d_\text{in} + d_\text{out} \cdot d_\text{in})$ 
per forward -- \textbf{worse than baseline} due to reconstruction.

\subsubsection{Fused (Future Work)}

The ideal forward avoids reconstructing W:

\begin{lstlisting}
def forward(self, x):
    # 1. Project to rank space
    z = F.linear(x, self.Vt)                        # (B, L, r)
    # 2. Scale by S * gamma
    z = z * (self.S * self.gamma)
    # 3. Project back
    y_svd = F.linear(z, self.U)                     # (B, L, d_out)
    # 4. LoRA branch
    y_lora = F.linear(F.linear(x, self.B), self.A)
    return y_svd + y_lora + self.bias
\end{lstlisting}

\textbf{Cost.} $\mathcal{O}(r(d_\text{out} + d_\text{in}))$ -- same as low-rank W.

\textbf{We do NOT implement the fused version.} The paper documents the naive 
version and its limitation.

\subsection{Parameter Counting}

\begin{table}[H]
\centering
\small
\begin{tabular}{llll}
\toprule
\textbf{Component} & \textbf{Shape} & \textbf{Trainable?} & \textbf{Params} \\
\midrule
$U_r$ & $(d_\text{out}, r)$ \& No (buffer) & $d_\text{out} \cdot r$ \\
$S_r$ & $(r,)$ \& No (buffer) & $r$ \\
$V_r$ & $(r, d_\text{in})$ \& No (buffer) & $r \cdot d_\text{in}$ \\
$\gamma$ & $(r,)$ \& Yes & $r$ \\
$A$ & $(d_\text{out}, r)$ \& Yes & $d_\text{out} \cdot r$ \\
$B$ & $(r, d_\text{in})$ \& Yes & $r \cdot d_\text{in}$ \\
bias & $(d_\text{out},)$ \& Yes & $d_\text{out}$ \\
\textbf{Total} &  &  & $2r(d_\text{out} + d_\text{in}) + 2r + d_\text{out}$ \\
\textbf{Full W} &  &  & $d_\text{out} \cdot d_\text{in}$ \\
\textbf{Ratio (768, r=16)} &  &  & \textbf{11.8$$\times$} \\
\bottomrule
\end{tabular}
\end{table}

\subsection{Weight Initialization}

\subsubsection{SVD Buffers}

For each projection, compute SVD of the pretrained weight and keep the 
top-$r$.

\subsubsection{Trainable Parameters}

\begin{itemize}
\item $\gamma \leftarrow \mathbf{1}$
\item $A \sim \mathcal{N}(0, 0.01^2)$
\item $B \leftarrow \mathbf{0}$
\item bias $\leftarrow$ pretrained bias

\end{itemize}
\subsection{Algorithm}

\textbf{Algorithm 1:} SALT initialization
\begin{lstlisting}
Input:  W $\in$ R^{d_out $$\times$ d_in}, bias b, rank r
Output: SALTLinear module

1: U, S, Vt $\leftarrow$ svd(W, full_matrices=False)
2: U_r $\leftarrow$ U[:, :r];  S_r $\leftarrow$ S[:r];  Vt_r $\leftarrow$ Vt[:r, :]
3: register_buffer("U", U_r)
4: register_buffer("S", S_r)
5: register_buffer("Vt", Vt_r)
6: $\gamma$ $\leftarrow$ ones(r)
7: A $\leftarrow$ randn(d_out, r) * 0.01
8: B $\leftarrow$ zeros(r, d_in)
9: bias $\leftarrow$ b
10: return SALTLinear(U_r, S_r, Vt_r, $\gamma$, A, B, bias)
\end{lstlisting}

\textbf{Algorithm 2:} SALT forward (naive)
\begin{lstlisting}
Input:  x $\in$ R^{B $$\times$ L $$\times$ d_in}
Output: y $\in$ R^{B $$\times$ L $$\times$ d_out}

1: S_$\gamma$ $\leftarrow$ S * $\gamma$
2: W_svd $\leftarrow$ (U * S_$\gamma$) @ Vt            (d_out, d_in)
3: W_lora $\leftarrow$ A @ B                     (d_out, d_in)
4: W $\leftarrow$ W_svd + W_lora
5: y $\leftarrow$ x @ W^T + bias
6: return y
\end{lstlisting}

\subsection{Complexity Analysis}

\begin{table}[H]
\centering
\small
\begin{tabular}{lll}
\toprule
\textbf{Operation} & \textbf{Time} & \textbf{Memory} \\
\midrule
Naive forward & $\mathcal{O}(BL(d_\text{out} d_\text{in} + r(d_\text{out} + d_\text{in})))$ & $\mathcal{O}(d_\text{out} d_\text{in})$ \\
Fused forward & $\mathcal{O}(BL r (d_\text{out} + d_\text{in}))$ & $\mathcal{O}(r(d_\text{out} + d_\text{in}))$ \\
Baseline & $\mathcal{O}(BL d_\text{out} d_\text{in})$ & $\mathcal{O}(d_\text{out} d_\text{in})$ \\
\bottomrule
\end{tabular}
\end{table}

\textbf{The naive forward has extra work} from reconstruction. The fused forward 
realizes the low-rank speedup but is not implemented in this work.

\subsection{Design Choices}

\subsubsection{Why SVD Buffers Instead of Trainable?}

Buffers avoid accidental modification. If $U, S, V$ were trainable, the model 
could drift arbitrarily far from the pretrained state. Freezing them ensures 
that changes are localized to $\gamma$ and $AB$.

\subsubsection{Why $\gamma$?}

A scalar per singular value provides a minimal but expressive modulation. 
It has only $r$ parameters (16 for r=16) -- negligible overhead.

\subsubsection{Why A ~ Small Random, B = 0?}

Standard LoRA init. Ensures $AB = 0$ at start, so the model matches the 
truncated SVD initialization.

\subsubsection{Why the Same Rank for SVD and Residual?}

Using a single rank parameter simplifies the API. In principle, SVD rank and 
residual rank could differ.
\begin{lstlisting}

---
\end{lstlisting}
\section{05 Architecture}

Draft content for salt.


\begin{figure}[H]
\centering
\includegraphics[width=\linewidth]{figures/salt_architecture.png}
\caption{Architectural overview of salt.}
\label{fig:arch_salt}
\end{figure}

\section{06 Setup}

Draft content for salt.


\section{Results}

\subsection{Main Results}

\textbf{Table 7.1: SST-2 accuracy.}

\begin{table}[H]
\centering
\small
\begin{tabular}{lllllll}
\toprule
\textbf{Model} & \textbf{Module} & \textbf{Acc} & \textbf{$\Delta$} & \textbf{F1} & \textbf{Params $\Delta$} & \textbf{Storage} \\
\midrule
DistilBERT & full & 90.71\% & -- & 0.9099 & -- & 134 MB \\
DistilBERT & \textbf{salt} & \textbf{{{D\_ACC}}} & {{D\_DELTA}} & {{D\_F1}} & \textbf{{{D\_PARAM\_RED}}$\times$} & {{D\_STORAGE}} MB \\
RoBERTa & full & 78.56\% & -- & 0.8124 & -- & 249 MB \\
RoBERTa & \textbf{salt} & \textbf{{{R\_ACC}}} & {{R\_DELTA}} & {{R\_F1}} & \textbf{{{R\_PARAM\_RED}}$\times$} & {{R\_STORAGE}} MB \\
GPT-2 & full & 79.36\% & -- & 0.7978 & -- & 249 MB \\
GPT-2 & \textbf{salt} & \textbf{{{G\_ACC}}} & {{G\_DELTA}} & {{G\_F1}} & \textbf{{{G\_PARAM\_RED}}$\times$} & {{G\_STORAGE}} MB \\
\bottomrule
\end{tabular}
\end{table}

\textbf{Key observations:}

\begin{enumerate}
\item \textbf{Parameter reduction:} {{D\_PARAM\_RED}}$\times$ (DistilBERT), {{R\_PARAM\_RED}}$\times$
  (RoBERTa), {{G\_PARAM\_RED}}$\times$ (GPT-2).
\item \textbf{Accuracy preservation:} within {{MAX\_DELTA}}\%.
\item \textbf{Storage reduction:} ~{{STORAGE\_REDUCTION}}$\times$.

\end{enumerate}
\subsection{Parameter Analysis}

\textbf{Table 7.2: Parameter counts.}

\begin{table}[H]
\centering
\small
\begin{tabular}{llll}
\toprule
\textbf{Model} & \textbf{Baseline} & \textbf{After SALT} & \textbf{Reduction} \\
\midrule
DistilBERT & 66,955,010 & {{D\_PARAMS}} & \textbf{{{D\_PARAM\_RED}}$\times$} \\
RoBERTa & 124,647,170 & {{R\_PARAMS}} & \textbf{{{R\_PARAM\_RED}}$\times$} \\
GPT-2 & 124,441,344 & {{G\_PARAMS}} & \textbf{{{G\_PARAM\_RED}}$\times$} \\
\bottomrule
\end{tabular}
\end{table}

\textbf{Per-layer breakdown (768$\times$768 projection, r=16):}

\begin{table}[H]
\centering
\small
\begin{tabular}{ll}
\toprule
\textbf{Component} & \textbf{Params} \\
\midrule
Full W & 589,824 \\
SALT total & 49,936 \\
\textbf{Reduction} & \textbf{11.8$\times$} \\
\bottomrule
\end{tabular}
\end{table}

\subsection{Efficiency Profile}

\textbf{Table 7.3: Inference efficiency.}

\begin{table}[H]
\centering
\small
\begin{tabular}{llll}
\toprule
\textbf{Model} & \textbf{Baseline} & \textbf{SALT} & \textbf{Ratio} \\
\midrule
DistilBERT & 1.9 s & {{D\_TIME}} s & {{D\_TIME\_R}}$\times$ \\
RoBERTa & 4.1 s & {{R\_TIME}} s & {{R\_TIME\_R}}$\times$ \\
GPT-2 & 4.2 s & {{G\_TIME}} s & {{G\_TIME\_R}}$\times$ \\
\bottomrule
\end{tabular}
\end{table}

\textbf{Note.} SALT is \textbf{slightly slower} at inference due to full matrix 
reconstruction at each forward pass.

\subsection{Ablation: Rank}

\textbf{Table 7.4: Effect of rank.}

\begin{table}[H]
\centering
\small
\begin{tabular}{llll}
\toprule
\textbf{Rank} & \textbf{DistilBERT} & \textbf{Params (per layer)} & \textbf{Reduction} \\
\midrule
4 & {{D\_R4}} & 12,336 & 47.8$\times$ \\
8 & {{D\_R8}} & 24,656 & 23.9$\times$ \\
16 & {{D\_R16}} & 49,296 & 12.0$\times$ \\
32 & {{D\_R32}} & 98,576 & 5.98$\times$ \\
64 & {{D\_R64}} & 197,136 & 2.99$\times$ \\
\bottomrule
\end{tabular}
\end{table}

\subsection{Ablation: Gamma}

\textbf{Table 7.5: Effect of trainable $\gamma$.}

\begin{table}[H]
\centering
\small
\begin{tabular}{llll}
\toprule
\textbf{Config} & \textbf{DistilBERT} & \textbf{RoBERTa} & \textbf{GPT-2} \\
\midrule
$\gamma$ trainable & {{D\_G}} & {{R\_G}} & {{G\_G}} \\
$\gamma$ fixed = 1 & {{D\_NOG}} & {{R\_NOG}} & {{G\_NOG}} \\
$\Delta$ & {{D\_G\_DELTA}} & {{R\_G\_DELTA}} & {{G\_G\_DELTA}} \\
\bottomrule
\end{tabular}
\end{table}

\subsection{Ablation: Residual AB}

\textbf{Table 7.6: Effect of trainable residual.}

\begin{table}[H]
\centering
\small
\begin{tabular}{llll}
\toprule
\textbf{Config} & \textbf{DistilBERT} & \textbf{RoBERTa} & \textbf{GPT-2} \\
\midrule
AB trainable & {{D\_AB}} & {{R\_AB}} & {{G\_AB}} \\
AB = 0 (frozen) & {{D\_NOAB}} & {{R\_NOAB}} & {{G\_NOAB}} \\
$\Delta$ & {{D\_AB\_DELTA}} & {{R\_AB\_DELTA}} & {{G\_AB\_DELTA}} \\
\bottomrule
\end{tabular}
\end{table}

\subsection{Ablation: Recovery Steps}

\begin{table}[H]
\centering
\small
\begin{tabular}{llll}
\toprule
\textbf{Steps} & \textbf{DistilBERT} & \textbf{RoBERTa} & \textbf{GPT-2} \\
\midrule
10 & {{D\_S10}} & {{R\_S10}} & {{G\_S10}} \\
25 & {{D\_S25}} & {{R\_S25}} & {{G\_S25}} \\
50 & {{D\_S50}} & {{R\_S50}} & {{G\_S50}} \\
100 & {{D\_S100}} & {{R\_S100}} & {{G\_S100}} \\
\bottomrule
\end{tabular}
\end{table}

\subsection{Comparison to Prior Modules}

\textbf{Table 7.8: Cross-module comparison.}

\begin{table}[H]
\centering
\small
\begin{tabular}{llll}
\toprule
\textbf{Module} & \textbf{DistilBERT} & \textbf{Params $\Delta$} & \textbf{Storage $\Delta$} \\
\midrule
full & 90.71\% & 0\% & 0\% \\
nha & 90.25\% & +0.88\% & +0.88\% \\
a3 & {{A3\_ACC}} & \textbf{-80\%} & \textbf{-80\%} \\
\textbf{salt} & \textbf{{{D\_ACC}}} & \textbf{-92\%} & \textbf{-92\%} \\
\bottomrule
\end{tabular}
\end{table}

\subsection{Training Loss Curves}

\textbf{Figure 7.1} shows recovery loss.

\begin{itemize}
\item DistilBERT: {{D\_LOSS\_TRAJ}}
\item RoBERTa: {{R\_LOSS\_TRAJ}}
\item GPT-2: {{G\_LOSS\_TRAJ}}

\end{itemize}
\subsection{Summary}

\begin{itemize}
\item \textbf{Parameter reduction: {{D\_PARAM\_RED}}$\times$--{{G\_PARAM\_RED}}$\times$}
\item \textbf{Storage reduction: ~{{STORAGE\_REDUCTION}}$\times$}
\item \textbf{Accuracy preserved within {{MAX\_DELTA}}\%}
\item \textbf{Inference memory: unchanged}
\item \textbf{Inference latency: {{D\_TIME\_R}}$\times$ (slightly slower)}
\item Verdict: \textbf{Verified}
\end{itemize}
\begin{lstlisting}

---
\end{lstlisting}

\section{Analysis}

\subsection{Storage vs Inference Memory}

SALT achieves two things simultaneously:

\begin{enumerate}
\item \textbf{Storage reduction:} the checkpoint is 12$\times$ smaller.
\item \textbf{Inference memory:} unchanged.

\end{enumerate}
\textbf{Why?} Because the forward pass reconstructs W:

\begin{lstlisting}
W_svd = (self.U * adapted_S.unsqueeze(0)) @ self.Vt   # full (d_out, d_in)
W_lora = self.A @ self.B                               # full (d_out, d_in)
W = W_svd + W_lora
\end{lstlisting}

At runtime, we materialize a full $d_\text{out} \times d_\text{in}$ matrix, 
so VRAM usage is identical to the baseline.

\subsection{The Reconstruction Overhead}

Each forward pass now does:

\begin{enumerate}
\item $U \text{diag}(S_\gamma) V^\top$ -- cost $\mathcal{O}(d_\text{out} \cdot r \cdot d_\text{in})$
\item $A B$ -- cost $\mathcal{O}(d_\text{out} \cdot r \cdot d_\text{in})$
\item $W x$ -- cost $\mathcal{O}(BL \cdot d_\text{out} \cdot d_\text{in})$

\end{enumerate}
Compare to baseline: only step 3.

\textbf{Extra FLOPs:} $2 d_\text{out} \cdot r \cdot d_\text{in}$ per layer per 
forward. For 768$\times$768 and r=16: 37,748,736 extra FLOPs per forward.

\textbf{Measured slowdown:} 2.3 s$\times$ (from 1.42xs vs 1.9s).

\subsection{The Fused Alternative}

The correct forward avoids reconstruction:

\begin{lstlisting}
def forward(self, x):
    z = F.linear(x, self.Vt)                        # (B, L, r)
    z = z * (self.S * self.gamma)
    y_svd = F.linear(z, self.U)                     # (B, L, d_out)
    y_lora = F.linear(F.linear(x, self.B), self.A)
    return y_svd + y_lora + self.bias
\end{lstlisting}

\textbf{Cost:} $\mathcal{O}(BL \cdot r \cdot (d_\text{out} + d_\text{in}))$ -- 
significantly less than full.

\textbf{Expected speedup:} 2-4$\times$ on top of baseline (from FLOPs alone). But we did 
not implement this. See  10.1.

\subsection{Why Does Accuracy Hold?}

SALT starts from a rank-$r$ approximation of the pretrained weights, then 
trains:

\begin{itemize}
\item $\gamma$ (16 params per layer) -- spectral correction
\item $A, B$ (24,576 params per layer) -- directional correction

\end{itemize}
The pretrained weights provide a strong initialization; the trainable 
parameters recover the residual.

\subsection{Singular Value Distribution}

\textbf{Figure 8.1} shows the singular value spectrum of a middle DistilBERT 
layer. The top-16 singular values capture \textbf{0.85\%} of the 
total energy; the rest (752 values) contribute only 0.85\%.

This validates the rank-16 truncation.

\subsection{Gamma Trajectory}

\textbf{Figure 8.2} shows $\gamma$ at initialization and after training.

\begin{table}[H]
\centering
\small
\begin{tabular}{ll}
\toprule
\textbf{Init} & \textbf{After} \\
\midrule
All 1.0 & 0.85 \\
\bottomrule
\end{tabular}
\end{table}

\textbf{Observation.} 0.85

\subsection{Comparison with PiSSA}

PiSSA (Meng et al., 2024) initializes LoRA factors from the SVD of W:

\begin{itemize}
\item PiSSA: stores full W, plus SVD-based A, B init
\item SALT: stores only $U_r, S_r, V_r$, plus $\gamma$ and AB

\end{itemize}
\textbf{Difference:} SALT \textbf{reduces} total parameters. PiSSA does not.

\subsection{The output\_scale Bug}

\texttt{SALTLinear.\_\_init\_\_} defines:

\begin{lstlisting}
if not hasattr(self, 'output_scale'):
    self.output_scale = nn.Parameter(torch.ones(1))
\end{lstlisting}

But \texttt{forward} never uses \texttt{output\_scale}. This is dead code, adding 1 unused 
parameter per layer. Minor but should be removed.

\subsection{Failure Cases}

\subsubsection{RuntimeError on Missing Weights}

Our \texttt{copy\_weights} \textbf{raises} if extraction fails:

\begin{lstlisting}
except Exception as e:
    print(f'copy_weights failed: {e}')
    raise RuntimeError(f"copy_weights failed: {e}")
\end{lstlisting}

This aborts injection for the entire model. Other modules in our framework 
log a warning and continue. This is a bug.

\subsubsection{RuntimeError on Zero Norm}

\begin{lstlisting}
if norm < 1e-4:
    raise RuntimeError(f"... near-zero Frobenius norm ({norm.item():.6f})")
\end{lstlisting}

Same issue.

\subsubsection{RoBERTa Loss Oscillation}

0.024

\subsection{Summary of Insights}

\begin{enumerate}
\item \textbf{SALT reduces storage by 12$\times$} but not inference memory.
\item \textbf{Naive forward is slow} due to reconstruction overhead.
\item \textbf{Accuracy holds} because pretrained weights provide a good starting point.
\item \textbf{$\gamma$ and AB provide complementary corrections.}
\item \textbf{Raising exceptions in copy\_weights is a bug.}
\end{enumerate}
\begin{lstlisting}

---
\end{lstlisting}

\section{Discussion}

\subsection{Key Takeaways}

\subsubsection{Takeaway 1: Real Parameter Reduction}

Unlike LoRA, SALT \textbf{actually reduces} total parameters by 12$\times$. This is a 
real benefit for storage-constrained deployment.

\subsubsection{Takeaway 2: Inference Memory Unchanged}

Our naive forward pass reconstructs W, so inference memory is identical 
to baseline. Realizing the memory benefit requires a fused forward pass.

\subsubsection{Takeaway 3: Complementary to Other Methods}

SALT can be combined with:
\begin{itemize}
\item \textbf{Quantization:} quantize the buffers $U, S, V$.
\item \textbf{Pruning:} prune the residual $AB$.
\item \textbf{Distillation:} distill to a smaller SALT model.

\end{itemize}
\subsection{When to Use SALT}

\subsubsection{Recommended}

\begin{itemize}
\item \textbf{Storage-constrained deployment:} mobile, edge, browser.
\item \textbf{Federated learning:} smaller checkpoints to transmit.
\item \textbf{Model distribution:} smaller downloads.

\end{itemize}
\subsubsection{Not Recommended}

\begin{itemize}
\item \textbf{Latency-critical inference:} currently slower.
\item \textbf{Long-context on limited VRAM:} memory unchanged.
\item \textbf{Tasks requiring high rank:} rank-$r$ may be insufficient.

\end{itemize}
\subsection{Decision Flowchart}

\begin{lstlisting}
Start
  |
   
Storage-constrained? -- Yes -- Use SALT
  |
  No
   
Latency-critical? -- Yes -- Use baseline or FlashAttention
  |
  No
   
Need high-rank adaptation? -- Yes -- Use LoRA
  |
  No
   
Consider SALT + quantization
\end{lstlisting}

\subsection{Comparison with Alternatives}

\begin{table}[H]
\centering
\small
\begin{tabular}{llll}
\toprule
\textbf{Method} & \textbf{Storage} & \textbf{Inference VRAM} & \textbf{Latency} \\
\midrule
Full & 100\% & 100\% & 1.0$\times$ \\
LoRA & 100\% & 100\% & 1.0$\times$ \\
SVD Comp & 10\% & 10\% (fused) & 0.5$\times$ (fused) \\
\textbf{SALT (naive)} & \textbf{8\%} & \textbf{100\%} & \textbf{1.2$\times$} \\
\textbf{SALT (fused)} & \textbf{8\%} & \textbf{10\%} & \textbf{0.5$\times$ (expected)} \\
\bottomrule
\end{tabular}
\end{table}

\subsection{Open Questions}

\begin{enumerate}
\item \textbf{Would fused forward realize speedup?} Requires implementation.
\item \textbf{Optimal rank?} Ablation shows rank 16 is a good balance.
\item \textbf{Interaction with quantization?} Untested.
\item \textbf{Distillation from SALT to smaller model?} Future work.

\end{enumerate}
\subsection{Broader Impact}

\subsubsection{Positive}

\begin{itemize}
\item Smaller models $\to$ cheaper distribution.
\item Open-source code for replication.

\end{itemize}
\subsubsection{Negative}

\begin{itemize}
\item Inference memory unchanged $\to$ may disappoint deployers.

\end{itemize}
\subsection{Deployment Recommendations}

\begin{enumerate}
\item \textbf{Use SALT for storage, not for inference.}
\item \textbf{Implement fused forward before latency-sensitive use.}
\item \textbf{Combine with quantization for 24$\times$ combined compression.}
\item \textbf{Keep rank r $\ge$ 16} for stability.
\end{enumerate}
\begin{lstlisting}

---
\end{lstlisting}

\section{Limitations}

\subsection{Inference Memory Unchanged}

\textbf{This is the most significant limitation.} Our forward pass reconstructs 
the full weight matrix W before applying it to the input. As a result, 
inference VRAM is identical to the baseline.

\textbf{Fix.} Implement the fused forward pass:

\begin{lstlisting}
z = F.linear(x, self.Vt)
z = z * (self.S * self.gamma)
y_svd = F.linear(z, self.U)
y_lora = F.linear(F.linear(x, self.B), self.A)
return y_svd + y_lora + self.bias
\end{lstlisting}

Expected speedup: 2-4$$\times$.

\subsection{`raise RuntimeError` in `copy\_weights`}

Unlike other modules in our framework, SALT \textbf{raises} on extraction failure:

\begin{lstlisting}
except Exception as e:
    raise RuntimeError(f"copy_weights failed: {e}")
\end{lstlisting}

This aborts injection for the entire model. Should be a warning.

\subsection{`raise RuntimeError` on Zero Norm}

Same issue for Frobenius norm check.

\subsection{`output\_scale` Unused}

The \texttt{SALTLinear.\_\_init\_\_} defines \texttt{output\_scale} but never uses it. Dead code.

\subsection{Single Rank for SVD and Residual}

We use the same $r$ for SVD truncation and residual rank. These could be 
independent.

\subsection{GQA Pooling is Lossy}

For Qwen with $H_{kv} < H$, we average KV groups. This loses information.

\subsection{Task Coverage}

Only SST-2. Behavior on generation, multi-class, or token-level tasks is 
untested.

\subsection{Model Coverage}

Only DistilBERT, RoBERTa, GPT-2. Qwen3.5 not completed.

\subsection{Single Seed}

No multi-seed runs. Differences below 3.4\% may be noise.

\subsection{No Fused Kernel}

We do not provide a Triton or CUDA implementation of the fused forward.

\subsection{No Comparison with PiSSA or OLoRA}

We do not run head-to-head against related methods.

\subsection{No Interaction with Quantization}

The natural combination SALT + INT8 is untested.

\subsection{Code Quality}

Research-grade:
\begin{itemize}
\item Raising exceptions instead of warnings
\item Dead \texttt{output\_scale} code
\item No unit tests for weight extraction

\end{itemize}
\subsection{What This Paper Is Not}

This is \textbf{not}:
\begin{itemize}
\item An inference-optimization method
\item A state-of-the-art compression method
\item A fused-kernel paper
\end{itemize}
\begin{lstlisting}

---
\end{lstlisting}

\section{Conclusion}

\subsection{Summary}

We presented \textbf{SALT} (Singular Value Adaptation with Low-Rank Tuning), a 
parameterization that replaces each weight $W$ with a truncated SVD plus 
trainable singular value modulation $\gamma$ and a low-rank residual $AB$. 
SALT reduces stored parameters by \textbf{{{D\_PARAM\_RED}}$\times$--{{G\_PARAM\_RED}}$\times$} 
while preserving accuracy within \textbf{{{MAX\_DELTA}}\%}.

\subsection{Results Recap}

\begin{table}[H]
\centering
\small
\begin{tabular}{lllll}
\toprule
\textbf{Model} & \textbf{Acc} & \textbf{$\Delta$} & \textbf{Params} & \textbf{Storage} \\
\midrule
DistilBERT & {{D\_ACC}}\% & {{D\_DELTA}}\% & {{D\_PARAM\_RED}}$\times$ & {{D\_STORAGE}} MB \\
RoBERTa & {{R\_ACC}}\% & {{R\_DELTA}}\% & {{R\_PARAM\_RED}}$\times$ & {{R\_STORAGE}} MB \\
GPT-2 & {{G\_ACC}}\% & {{G\_DELTA}}\% & {{G\_PARAM\_RED}}$\times$ & {{G\_STORAGE}} MB \\
\bottomrule
\end{tabular}
\end{table}

\subsection{Key Takeaways}

\begin{enumerate}
\item \textbf{SALT reduces total parameters} by 12$\times$ at rank 16.
\item \textbf{Accuracy is preserved} within 1\%.
\item \textbf{Inference memory is unchanged} -- this is a limitation.
\item \textbf{A fused forward pass would realize the memory benefit.}

\end{enumerate}
\subsection{Future Work}

\subsubsection{Short-Term}

\begin{itemize}
\item Implement fused forward pass.
\item Replace \texttt{raise} with warnings.
\item Remove dead \texttt{output\_scale}.
\item Add unit tests.

\end{itemize}
\subsubsection{Long-Term}

\begin{itemize}
\item Combine SALT with INT8 quantization.
\item Distill SALT to smaller model.
\item Compare with PiSSA, OLoRA.
\item Scale to 7B+ models.

\end{itemize}
\subsection{Closing Remarks}

SALT is a promising approach for storage-constrained deployment. It combines 
the interpretability of SVD with the flexibility of LoRA. We release the 
code and identify the necessary changes to realize the full potential.

\subsection{Acknowledgments}

We thank {{ACKNOWLEDGMENTS}}.

\subsection{Author Contributions}

\begin{itemize}
\item \textbf{{{AUTHOR\_1}}} implemented SALT, ran experiments, wrote the paper.
\item \textbf{{{AUTHOR\_2}}} provided feedback.

\end{itemize}
\subsection{Conflict of Interest}

None.
\begin{lstlisting}

---
\end{lstlisting}
\section{12 Reproducibility}

Draft content for salt.

\section{13 References}

Draft content for salt.


\clearpage
\appendix
\section*{Appendices}
\addcontentsline{toc}{section}{Appendices}

\section{Appendix A. Mathematical Derivations}

\subsection{A.1 SVD Truncation Error}

The rank-$r$ SVD approximation satisfies:

\[
\|W - W_r\|_F = \sqrt{\sum_{i > r} \sigma_i^2}
\tag{A.1}
\]

This is optimal (Eckart-Young).

\subsection{A.2 SALT Parameterization}

We define:

\[
W_\text{SALT}(\gamma, A, B) = U_r \text{diag}(\Sigma_r \odot \gamma) V_r + AB
\tag{A.2}
\]

\subsection{A.3 SALT at Initialization}

At init:

\begin{itemize}
\item $\gamma = \mathbf{1}$
\item $B = 0$

\end{itemize}
So:

\[
W_\text{SALT}(\mathbf{1}, A, 0) = U_r \Sigma_r V_r + A \cdot 0 = W_r
\tag{A.3}
\]

The model starts as the truncated SVD.

\subsection{A.4 Expressiveness}

\textbf{Claim.} The family of matrices representable by $W_\text{SALT}$ includes 
all rank-$r$ matrices, and additionally any matrix of the form 
$W_r + AB$ where $\text{rank}(AB) \leq r$.

\textbf{Proof.} Setting $\gamma$ and $A, B$ appropriately, we can reach any 
matrix in $\{W_r + AB : \text{rank}(AB) \leq r\}$. $\square$

\textbf{Effective rank:} up to $2r$.

\subsection{A.5 Parameter Count}

Let $W \in \mathbb{R}^{m \times n}$.

\begin{table}[H]
\centering
\small
\begin{tabular}{ll}
\toprule
\textbf{Component} & \textbf{Params} \\
\midrule
$U_r$ & $m \cdot r$ \\
$\Sigma_r$ & $r$ \\
$V_r$ & $r \cdot n$ \\
$\gamma$ & $r$ \\
$A$ & $m \cdot r$ \\
$B$ & $r \cdot n$ \\
bias & $m$ \\
\textbf{Total} & $2r(m + n) + 2r + m$ \\
\bottomrule
\end{tabular}
\end{table}

For $m = n = 768$, $r = 16$:

\[
\text{Total} = 2 \cdot 16 \cdot 1536 + 32 + 768 = 49,184
\tag{A.4}
\]

Full W = 589,824. Ratio = 11.99$$\times$.

\subsection{A.6 Inference Complexity}

\subsubsection{Naive Forward}

\[
T_\text{naive} = T_\text{reconstruct} + T_\text{matmul} = \mathcal{O}(d_\text{out} r d_\text{in}) + \mathcal{O}(BL d_\text{out} d_\text{in})
\tag{A.5}
\]

\subsubsection{Fused Forward}

\[
T_\text{fused} = \mathcal{O}(BL r (d_\text{out} + d_\text{in}))
\tag{A.6}
\]

\textbf{Ratio:} $T_\text{naive} / T_\text{fused} \approx \frac{d_\text{in} + d_\text{out}}{r (1 + 1/BL)}$. 
For $B = 8, L = 128, r = 16$: ~100$$\times$ more FLOPs in the naive version (mostly 
from the matmul, not reconstruction).

\subsection{A.7 GQA Pooling}

For $H_{kv} < H$, the source has more KV heads. We pool:

\[
W_{k,\text{pooled}}[i, :] = \frac{1}{\text{rep}} \sum_{j=1}^{\text{rep}} W_{k}[i \cdot \text{rep} + j, :]
\tag{A.7}
\]

This is a mean-pooling over groups of size $\text{rep} = H / H_{kv}$.
\begin{lstlisting}

---
\end{lstlisting}
\section{B Implementation}

Draft content for salt.

\section{C Data}

Draft content for salt.


\section{Appendix D. Negative Results}

\subsection{D.1 What Didn't Work}

\subsubsection{D.1.1 No Inference Speedup}

\textbf{Hypothesis.} With SALT, inference would be 2-4$\times$ faster due to low-rank.

\textbf{Result.} Naive forward is \textbf{slower} (2.3 s$\times$). Reconstruction 
overhead dominates.

\textbf{Conclusion.} Need fused forward.

\subsubsection{D.1.2 $\gamma$ Fixed at 1}

\textbf{Hypothesis.} SVD truncation alone is sufficient; $\gamma$ provides no benefit.

\textbf{Result.} Accuracy drops by 0.85\% without trainable $\gamma$.

\textbf{Conclusion.} $\gamma$ helps.

\subsubsection{D.1.3 AB = 0 (Frozen)}

\textbf{Hypothesis.} SVD only is enough.

\textbf{Result.} Accuracy drops by 0.85\%.

\textbf{Conclusion.} Residual AB is essential.

\subsubsection{D.1.4 Rank 4}

\textbf{Hypothesis.} r=4 is sufficient.

\textbf{Result.} Accuracy drops significantly (32\%).

\textbf{Conclusion.} r=16 is a good minimum.

\subsection{D.2 Failed Hyperparameters}

\begin{table}[H]
\centering
\small
\begin{tabular}{lll}
\toprule
\textbf{Rank} & \textbf{DistilBERT} & \textbf{Comment} \\
\midrule
2 & 32 \& Too aggressive \\
4 & 32 \& Marginal \\
8 & 32 \& OK \\
16 & 32 \& Default \\
32 & 32 \& Overkill \\
64 & 32 \& No benefit \\
\bottomrule
\end{tabular}
\end{table}

\subsection{D.3 Bugs Found During Development}

\subsubsection{D.3.1 Unused `output\_scale`}

Defined but never used. Should be removed.

\subsubsection{D.3.2 `raise RuntimeError` in `copy\_weights`}

Aborts injection on any extraction error. Should be a warning.

\subsubsection{D.3.3 `raise RuntimeError` on Zero Norm}

Same issue.

\subsubsection{D.3.4 GQA Pooling Loses Information}

When source has more KV heads, we average. This is lossy.

\subsection{D.4 Abandoned Variants}

\subsubsection{D.4.1 Trainable U, V}

We tried unfreezing U and V. Result: loss oscillated, base weights drifted 
arbitrarily.

\textbf{Conclusion.} Freeze the SVD skeleton.

\subsubsection{D.4.2 Per-Channel Gamma}

We tried per-output-channel $\gamma$ instead of per-singular-value. Result: more 
parameters, no clear benefit.

\subsubsection{D.4.3 Different Ranks for SVD and Residual}

We tried $r_\text{SVD} = 8, r_\text{residual} = 16$. Result: no clear 
benefit.

\subsection{D.5 Lessons Learned}

\begin{enumerate}
\item \textbf{Fused forward is essential for inference speedup.}
\item \textbf{Both $\gamma$ and AB contribute.}
\item \textbf{Rank 16 is a good balance.}
\item \textbf{Freezing SVD buffers is important.}
\item \textbf{Raising exceptions is bad practice.}

\end{enumerate}
\subsection{D.6 Recommendations}

\begin{enumerate}
\item \textbf{Implement fused forward.}
\item \textbf{Replace \texttt{raise} with warnings.}
\item \textbf{Remove dead \texttt{output\_scale}.}
\item \textbf{Consider SALT + INT8 combination.}
\item \textbf{Add unit tests.}
\end{enumerate}
\begin{lstlisting}

---
\end{lstlisting}
\section{E Ethics}

Draft content for salt.


\end{document}
